%% ============================================================
\section{Conclusion}
\label{sec:conclusion}

We have presented SqC, a static analysis tool implementing 311 CERT C
rules across 17 categories.  Through tree-sitter-based parsing,
control-flow graph construction, null-state dataflow analysis,
value-range analysis, and inter-procedural function summaries
(including cross-function taint propagation), SqC achieves an 87.1\%
true positive rate on the Juliet benchmark with 43 CWEs at 100\%
precision, while flagging 12.9\% of the suite's planted flaw
lines---detection breadth, not output cleanliness, being the limit
that remains.  On a five-tool comparison, SqC surpasses Frama-C
(79.5\% vs.\ 61.0\% on overlapping CWEs) and now beats clang-tidy
outright on three of those CWEs, while covering 5--12$\times$
more vulnerability categories.

On real-world codebases totaling 1.12 million lines of code across
seven projects, total violations on the original five were reduced
82.4\% through iterative refinement.  SqC detects vulnerability
classes that existing open-source tools do not cover, including
POSIX thread misuse, integer overflow, and inter-procedural null
dereference.

The tool's primary limitation---its tree-sitter foundation---is also
its key advantage: SqC analyzes any C file without a build system,
compilation database, or preprocessing step.  For teams that need
broad CERT C coverage with minimal integration effort, this tradeoff
is worthwhile.

Much of the post-v0.4.22 progress came not from the alias analysis or
symbolic execution this paper once projected as necessary, but from
reframing rules as opt-in provenance gates (\cref{sec:provenance-frontier}):
a cheaper redesign than either, and one motivated by real-world
precision rather than Juliet score, that happened to lift the Juliet
rate as a side effect.  This suggests future work should continue
looking for rule-level reframings before reaching for the heavier
architectural investments.  That said, some residual gaps do appear to
require deeper analysis: the MEM31-C real-world false-positive rate
remains largely unmoved by a first ownership-tracking pass
(\cref{sec:limitations}) because its dominant cause---parameter-owned
vs.\ locally-owned struct lifetimes---is exactly the kind of question
a real ownership or alias model answers and a same-function heuristic
cannot.  Future work therefore includes: a parameter-ownership model
for MEM31-C; a broader audit of which other rules would benefit from
the registry-based macro-expansion engine now used by MEM30/31-C and
EXP33-C; scoring the existing per-rule fixture corpus as a third
benchmark tier, since 40\% of SqC's rules have no true-positive
evidence on either current benchmark and no further adjudication of a
mature-code corpus can supply it (\cref{sec:corpus-maturity}); alias
analysis for pointer-related rules generally; and
further expanding the cross-function taint infrastructure --- CWE-78
(OS command injection) now achieves 100\% precision via caller-aware
ENV03-C suppression and return-value propagation, and the same summary
bits extend naturally to CWE-89 sinks (\texttt{sqlite3\_exec},
\texttt{mysql\_query}, \texttt{PQexec}) and format-string injection.
